\subsubsection{Kernel Ridge Regression}

Trong phần này, ta triển khai Kernel Ridge Regression với hai loại kernel là
RBF và Polynomial, đồng thời so sánh với Ridge Regression tuyến tính.

Đối với RBF kernel, tham số $\gamma$ được lựa chọn thông qua grid search với
tập giá trị $\{0.01, 0.1, 1\}$. Kết quả tốt nhất đạt được tại $\gamma = 0.1$.

Đối với Polynomial kernel, các tham số được tìm kiếm trên tập giá trị degree
là $\{2, 3\}$ và coef0 là $\{0, 1\}$ Kết quả tối ưu là degree = 2 và coef0
= 1.

\tablename~\ref{tab:kernel-ridge-cross-validation} trình bày kết quả
cross-validation của các mô hình. RBF Ridge cho kết quả tốt nhất, trong khi
Ridge tuyến tính kém nhất và Polynomial Ridge có độ biến động lớn.

\paragraph{Kết quả cross-validation}

\begin{table}[htbp]
  \caption{Kết quả cross-validation của các mô hình}
  \label{tab:kernel-ridge-cross-validation}
  \centering
  \import{\tablespath}{kernel-ridge-cross-validation}
\end{table}

\paragraph{Kết quả trên tập kiểm tra}
\tablename~\ref{tab:kernel-ridge-test-result} trình bày kết quả trên tập kiểm
tra của các mô hình. Có thể thấy rằng mô hình RBF Ridge tiếp tục đạt hiệu năng
tốt nhất với các giá trị MSE, RMSE và MAE thấp nhất, đồng thời hệ số $R^2$ cao
nhất, cho thấy khả năng tổng quát hóa tốt. Polynomial Ridge cũng cải thiện so
với Ridge tuyến tính, tuy nhiên vẫn kém hơn RBF Ridge. Trong khi đó, Ridge
Regression tuyến tính cho kết quả thấp nhất do không thể mô hình hóa các quan
hệ phi tuyến trong dữ liệu.

\begin{table}[htbp]
  \caption{Kết quả trên tập kiểm tra}
  \label{tab:kernel-ridge-test-result}
  \centering
  \import{\tablespath}{kernel-ridge-test-result}
\end{table}

\paragraph{Kiểm định thống kê}

\tablename~\ref{tab:kernel-ridge-statistical-test} cho thấy kết quả kiểm định
thống kê giữa các mô hình. Cụ thể, sự khác biệt giữa Ridge và RBF Ridge là có
ý nghĩa thống kê với p-value rất nhỏ trong cả t-test ($1.99 \times 10^{-13}$)
và Wilcoxon test (0.001953), nhỏ hơn mức ý nghĩa $\alpha = 0.05$. Điều này xác
nhận rằng RBF Ridge cải thiện hiệu năng một cách đáng kể so với Ridge tuyến
tính. Đối với cặp Ridge và Polynomial Ridge, cả hai kiểm định đều cho p-value
lớn hơn 0.05 (0.5153 và 0.4316), cho thấy không có sự khác biệt có ý nghĩa
thống kê giữa hai mô hình này. Điều này phù hợp với các kết quả thực nghiệm
trước đó, khi Polynomial Ridge không cho thấy sự cải thiện ổn định so với
Ridge. Đối với cặp RBF Ridge và Polynomial Ridge, kết quả không hoàn toàn nhất
quán giữa hai kiểm định: t-test cho p-value = 0.2299 (không có ý nghĩa thống
kê), trong khi Wilcoxon test cho p-value = 0.0020 (có ý nghĩa thống kê). Sự
khác biệt này có thể xuất phát từ việc t-test giả định phân phối chuẩn, trong
khi Wilcoxon test là kiểm định phi tham số và nhạy hơn với sự khác biệt về phân
phối. Tổng hợp lại, có thể kết luận rằng RBF Ridge là mô hình vượt trội nhất,
với sự cải thiện có ý nghĩa thống kê so với Ridge, trong khi Polynomial Ridge
không cho thấy sự khác biệt rõ ràng và ổn định so với các mô hình còn lại.

\begin{table}[htbp]
  \caption{So sánh thống kê giữa các mô hình}
  \label{tab:kernel-ridge-statistical-test}
  \centering
  \import{\tablespath}{kernel-ridge-statistical-test}
\end{table}

\paragraph{Nhận xét}

Kết quả cho thấy mô hình RBF Ridge đạt hiệu năng tốt nhất trên cả cross-
validation và tập kiểm tra, với sai số thấp nhất và hệ số $R^2$ cao nhất. Điều
này cho thấy dữ liệu có mối quan hệ phi tuyến mà kernel RBF có thể nắm bắt
hiệu quả.

Polynomial Ridge cải thiện kết quả so với Ridge tuyến tính trên tập kiểm tra,
tuy nhiên kết quả cross-validation có độ biến động lớn, cho thấy mô hình không
ổn định và có khả năng overfitting.

Kiểm định thống kê cho thấy sự khác biệt giữa Ridge và RBF Ridge là có ý nghĩa
thống kê (p-value rất nhỏ), trong khi sự khác biệt giữa Ridge và Polynomial
Ridge không có ý nghĩa thống kê (p-value lớn hơn 0.05).

Tóm lại, RBF kernel là lựa chọn phù hợp nhất trong bài toán này, trong khi
Ridge tuyến tính đóng vai trò baseline và Polynomial kernel cần được điều
chỉnh cẩn thận để đạt hiệu quả tốt.
